\documentclass[11pt]{article}
\usepackage{fontspec}
\setmainfont{Times New Roman}
\setsansfont{Arial}
\setmonofont{Consolas}
\usepackage{amsmath,amssymb}
\usepackage{geometry}
\usepackage{microtype}
\geometry{a4paper,margin=3cm}
\setlength{\parindent}{1.25em}
\setlength{\parskip}{0pt}
\emergencystretch=30em
\allowdisplaybreaks
\raggedbottom
\providecommand{\mM}{\mathfrak M}
\providecommand{\mG}{\mathfrak G}
\providecommand{\mU}{\mathfrak U}
\providecommand{\mB}{\mathfrak B}
\providecommand{\ma}{\mathfrak a}
\providecommand{\mb}{\mathfrak b}
\providecommand{\me}{\mathfrak e}
\providecommand{\mx}{\mathfrak x}

\begin{document}

% Унит: Папер 17 Сецтион 07. Дирецт Интерславиц Латин транслатион фром те Герман финал-аудитед соурце слице.

\setcounter{footnote}{15}
\subsection*{\S~7. Случај једнозначного разклада.}

Неха
\[
 \mG=\mU_1+\cdots+\mU_k=\mB_1+\cdots+\mB_l
\]
буде група остатков, јеј составне чести \(\mU_1,\ldots,\mU_k\), \(\mB_1,\ldots,\mB_l\) сут иредуцибилне. Тогда имајемо
\begin{equation}
 \mx\me=\mx\ma_1+\cdots+\mx\ma_k
 =\mx\mb_1+\cdots+\mx\mb_l,
\tag{22}
\end{equation}
при чем \(\ma_\chi\) јест клас остатков в \(\mG\), изоморфно приредјены јединичному класу од \(\mU_\chi\), а \(\mb_\lambda\) јест клас остатков в \(\mG\), изоморфно приредјены јединичному класу од \(\mB_\lambda\). Ако ныне заменимо \(\mx\) с \(\mx\mb_\lambda\), кде \(\lambda\) јест фиксованы индекс из реда \(1,\ldots,l\), доставајемо
\begin{equation}
 \mx\mb_\lambda=\mx\mb_\lambda\ma_1+\cdots+\mx\mb_\lambda\ma_k.
\tag{23}
\end{equation}
То представјенье групы \(\mB_\lambda\) јест једнозначно, але не јест адитивны разклад од \(\mB_\lambda\), бо поједине класы остатков на право не мусет належати к \(\mB_\lambda\). Понеже \(\mb_\lambda\) не изчезаје, не всакы продукт \(\mb_\lambda\ma_\chi\) може быти нуловы. Без ограниченьа обчности теды неха
\[
 \mb_\lambda\ma_1\ne0,
 \ldots,
 \mb_\lambda\ma_\rho\ne0,
 \qquad
 \mb_\lambda\ma_{\rho+1}=0,
 \ldots,
 \mb_\lambda\ma_k=0.
\]
Тогда
\[
 \mx\mb_\lambda=\mx\mb_\lambda\ma_1+\cdots+\mx\mb_\lambda\ma_\rho.
\]

Ныне за фиксовано \(\chi\) из реда \(1,\ldots,\rho\) разматрјамо целокупност оных класов остатков \(\mx\mb_\lambda\), кторе не сут \(0\) и за кторе такоже \(\mx\mb_\lambda\ma_\chi\ne0\). Ты две системы, \(\mx\mb_\lambda\) и \(\mx\mb_\lambda\ma_\chi\), сут изоморфне. Бо прво, приредјенье јест једно-једнозначно: из \(\mx\mb_\lambda\ma_\chi=0\) по предположеньу следује \(\mx\mb_\lambda=0\), а из последньего заради једнозначности представјеньа (23) следује \(\mx\mb_\lambda\ma_\chi=0\). Далье, суме одповедаје сума, а продукту с либовольным полиномом одповедны продукт. При фиксованом \(\chi\) можемо ныне провести исто разсматрјанье за всако \(\lambda\), за кторо \(\mb_\lambda\ma_\chi\ne0\). К всакому \(\ma_\chi\) мора быти најмене једно тако \(\mb_\lambda\), понеже
\[
 \ma_\chi=(\mb_1+\cdots+\mb_l)\ma_\chi\ne0.
\]

Најпрво предположимо, же за всако \(\chi\) обстаје само једно \(\lambda\), и за всако \(\lambda\) само једно \(\chi\), тако же \(\mb_\lambda\ma_\chi\ne0\).

Из того приредјеньа следује \(k=l\). Далье можно избрати означеньа тако, же \(\mb_\chi\ma_\chi\ne0\) и \(\mb_\lambda\ma_\chi=0\) при \(\lambda\ne\chi\). Тогда кажемо, же продукты \(\mb_\lambda\ma_\chi\) творе диагоналну схему. В том случају по (23) јединица \(\mb_\chi=\mb_\chi\ma_\chi\), а заради
\begin{equation}
 \ma_\chi=\mb_1\ma_\chi+\cdots+\mb_l\ma_\chi=\mb_\chi\ma_\chi=\mb_\chi
\tag{24}
\end{equation}
обче \(\mU_\chi=\mB_\chi\). Теды обстаје једнозначност разклада, и продукты \(\ma_\chi\mb_\lambda\) такоже творе диагоналну схему.

Ако специјално комутативны закон важи за класы остатков \(\ma_\chi\mb_\lambda\), тогда горно предположенье јест всегда исполньено. Бо в (23) всака составна чест \(\mx\mb_\lambda\ma_\chi=\mx\ma_\chi\mb_\lambda\) тогда садржи се в \(\mB_\lambda\), тако же (23) даваје адитивны разклад групы \(\mB_\lambda\). Заради иредуцибилности \(\mB_\lambda\) всако \(\mb_\lambda\ma_\chi\), кроме једного, јест \(=0\). Заради (24) јешче при фиксованом \(\chi\) само једно \(\mb_\lambda\ma_\chi\) јест различно од нулы, бо иначе (24), заради заменльивости \(\ma_\chi\) и \(\mb_\lambda\), дало бы адитивны разклад од \(\mU_\chi\), ктора мела быти иредуцибилна.\footnote{Одповедно важи в некомутативном двустранном случају, обработаном у \emph{Schmeidler} в цитованом месту, кде продукт всакого класа остатков једноје подгрупы с всакым класом остатков другоје подгрупы изчезаје. Тогда наиме заради
\[
\mb_\lambda\ma_\chi=\mb_\lambda\ma_\chi\mb_1+\cdots+\mb_\lambda\ma_\chi\mb_l
\]
и заради \(\mb_\lambda(\ma_\chi\mb_\mu)=0\) при \(\lambda\ne\mu\) такоже \(\mb_\lambda\ma_\chi=\mb_\lambda\ma_\chi\mb_\lambda\), теды јест садржен в \(\mB_\lambda\). Зато из иредуцибилности \(\mB_\lambda\) следује, же само за једно \(\chi\), точно за једно \(\chi\), продукт \(\mb_\lambda\ma_\chi\ne0\). В том самом способу показује се, же к всакому \(\ma_\chi\) обстаје једно и само једно тако \(\mb_\lambda\), же \(\ma_\chi\mb_\lambda\ne0\), из чего следује \(k=l\). Зато \(\mb_\lambda\ma_\chi\) творе диагоналну схему, чем доказана јест там доказана једнозначност.}

Ты саме закльучки важе в случају, же не за все, але само за чест \(\mb_\lambda\), \(\lambda=1,\ldots,\sigma\), исполньене сут условја \(\mb_\lambda\ma_\lambda\ne0\), \(\mb_\lambda\ma_\chi=0\) за \(\lambda\ne\chi\), \(\chi=1,\ldots,k\), и такоже за всако \(\mu=1,\ldots,l\), \(\mb_\mu\ma_\chi=0\), \(\mb_\mu\ma_\mu\ne0\), \(\chi=1,\ldots,\sigma\). Находимо \(\mU_\chi=\mB_\chi\) за \(\chi=1,\ldots,\sigma\); и ако положимо
\[
 \mU_{\sigma+1}+\cdots+\mU_k=\mU,
 \qquad
 \mB_{\sigma+1}+\cdots+\mB_l=\mB,
\]
тогда и за јединице \(\ma\) и \(\mb\) од \(\mU\) и \(\mB\) исполньено јест условје \(\mb\ma_\chi=0\), \(\mb_\chi\ma=0\) за \(\chi=1,\ldots,\sigma\), \(\mb\ma\ne0\), из чего следује \(\mU=\mB\).

\end{document}

